\subsection{A graph obstruction from a local contracted source algebra}

\begin{lemma}[Proved: a bounded local two-term complex]
\label{CP2:A-local-two-term-derived-amplitude}
Let $A$ be a finite-dimensional rational algebra with a unital character
$\chi:A\to\mathbb Q$ whose kernel $J$ is nilpotent.  Put
$k={}_\chi\mathbb Q$.  If $D$ is a complex of finite-dimensional right
$A$-modules concentrated in degrees zero and one, and
$D\otimes_A^{\mathbf L}k$ has bounded homology, then
\[
 H_j(D\otimes_A^{\mathbf L}k)=0\qquad(j\ge2).
\]
\end{lemma}
\begin{proof}
The nilpotence of $J$ implies that every element outside $J$ is a unit:
write it as a nonzero scalar times $1+z$ and invert $1+z$ by its finite
geometric series.  Hence $A$ is local, with radical $J$ and residue field
$\mathbb Q$.  A bounded-below free resolution of the finite complex $D$
can be chosen with finite free rank in each degree.  Cancel a pair of
adjacent rank-one free summands whenever a differential matrix has a
unit entry, using elementary changes of free bases.  Starting in the
lowest degree and proceeding upwards gives a minimal free resolution
$P\to D$: each differential matrix has entries in $J$.
There are only finitely many cancellations involving each fixed degree.

The complex $P\otimes_A k$ has zero differentials.  Boundedness of its
homology therefore says that $P_j\otimes_A k=0$ in all sufficiently
large degrees.  Since each $P_j$ is finite free, $P_j=0$ there.
Suppose its highest nonzero degree is $t\ge2$.  Then $H_t(P)=H_t(D)=0$,
so $d_t$ is injective.  Choose a nonzero $x\in A$ annihilated on the
left by $J$: if $J\ne0$, take $x$ in its last nonzero power, and if
$J=0$, take $x=1$.  For a free basis element $e_i$ of $P_t$, write
\[
 d_t(e_i)=\sum_j e_j a_{ji},\qquad a_{ji}\in J.
\]
Right-module linearity gives
$d_t(e_i x)=\sum_j e_j a_{ji}x=0$, while $e_i x\ne0$.
This contradicts injectivity.  Thus $P$ is concentrated in degrees at
most one, proving the claim.  No self-injectivity assumption on $A$ is
used.
\end{proof}

\begin{theorem}[Proved: graph chains over a projective or local row character]
\label{CP2:A-projective-local-row-graph-obstruction}
Let $M$ have a $2\times2$ rectangular minimal band $K$, put
$R=\mathbb QM$ and $A=R/\mathbb QK$, and let $V$ be its one-dimensional
left row-augmentation module.  Let a finite biset $S$ be left-$N$,
right-$M$, where $BN$ is rationally acyclic and the augmented chains of
$B(N,S)$ have a natural right-$M$-equivariant comparison to a finite graph
with strict edge maps, reversals and collapses as in
Lemma~\ref{CP2:A-six-graph-chain-obstruction}.
Suppose either $V$ is projective over $A$, or the kernel of its character
$\chi:A\to\mathbb Q$ is nilpotent.
If the collage has bounded rational homology, all its rational
degree-two products vanish.  There is no restriction on $|M|$ or $|S|$.
\end{theorem}
\begin{proof}
Collapse the actual corner $BN$.  This gives a rational homology
equivalence, and therefore a rational ring isomorphism, to
$Y=\operatorname{hocofib}(B(N,S,M)\to BM)$.
Use the augmented graph complex
\[
 C=[C_1(\Gamma;\mathbb Q)\longrightarrow
           \ker(C_0(\Gamma;\mathbb Q)\to\mathbb Q)]
\]
in degrees one and zero.  The reduced chains of $Y$ are computed by
$(C\otimes_R^{\mathbf L}\mathbb Q)[1]$ through the given equivariant
augmented comparison.

Write $e=x_{00}$ and $e_1=x_{01}$.  If $\Gamma e$ and $\Gamma e_1$
share a vertex, that vertex is fixed by every band actor.  It is then
fixed by every $m\in M$, since $em\in K$.  Its algebraic augmentation
section, transported by the natural chain comparison, makes the actual
bridge pullback injective on $H^2(BM;\mathbb Q)$.  The cofiber sequence
then makes every second class of $Y$ relative to $BM$.  These classes
lift from the actual suspension quotient and their products are zero.

Otherwise the two image graphs have disjoint vertices and edges.  The
literal multiplication map $Ce\otimes_{\mathbb Q}eR\to C$ is injective
degree by degree, with image $C\mathbb QK$: its two column copies have
disjoint supports, and in degree zero each has augmentation zero.
Its domain is a two-term complex of projective right $R$-modules.
Consequently the quotient
$D=C/(Ce\otimes_{\mathbb Q}eR)$ is a two-term right $A$-complex, and
removing the projective subcomplex affects no derived homology in
degrees at least three.  The exact sequence
$0\to V\to Re\to\mathbb Q\to0$, together with $De=0$, adds one
further shift.  Finally, the left-projective resolution
$0\to\mathbb QK\to R\to A\to0$ and the factorization
$\mathbb QK=Re\otimes_{\mathbb Q}eR$ show that no higher
change-of-rings term remains.  Thus, exactly as in the graph-chain
localization calculation,
\[
 H_n(Y;\mathbb Q)=H_{n-2}(D\otimes_A^{\mathbf L}V),\qquad n\ge4.
\]
These steps use the two-row, two-column band and require no bound on
the number of outside elements of $M$.

If $V$ is projective, the tensor on the right is represented by a
two-term complex and has no second homology.  In the other case,
boundedness of $Y$ gives boundedness of $D\otimes_A^{\mathbf L}V$;
the preceding lemma again kills its second homology.  Hence $H_4(Y)=0$
in either case, which kills every degree-two product.
\end{proof}

\begin{remark*}[Exact finite witnesses before applying the gate]
For a literal source table, the row character on an outside element
$a$ is the difference $\sigma_a(1)-\sigma_a(0)$ of its two row labels.
The contracted product of two outside basis elements is zero exactly
when their product lies in $K$.  Projectivity is certified by a rational
vector $p\in A$ with $ap=\chi(a)p$ on every basis element and
$\chi(p)=1$; these identities imply $p^2=p$ and $Ap\cong V$.
Nilpotence is certified by a finite sequence of rational subspaces
starting with $\ker\chi$, containing each successive right product by
the generators of $\ker\chi$, and ending at zero.  Every source-table
link and every displayed coefficient identity must be checked before
omitting its action scopes.  A source failing both sufficient tests is
retained in the finite search.
\end{remark*}
